\chapter[Lapisan Aplikasi: Komponen, Layanan, dan Integrasi]{Lapisan Aplikasi:
Komponen, Layanan, dan Integrasi Sistem Transfer Dana}
\label{ch:application-architecture}
\setlength{\emergencystretch}{2em}

\section*{Tujuan Pembelajaran}
\begin{itemize}
  \item Menjelaskan elemen lapisan aplikasi ArchiMate---komponen, layanan,
        fungsi, dan antarmuka aplikasi---serta relasi yang menautkannya ke
        proses bisnis dan objek data.
  \item Menyusun \emph{application landscape} untuk kapabilitas sasaran dan
        memetakan komponen ke layanan bisnis yang diwujudkannya.
  \item Menempatkan analisis aplikasi dalam TOGAF ADM \textbf{Fase C (Arsitektur
        Sistem Informasi)}, khususnya sub-fase \textbf{Arsitektur Aplikasi}:
        menyusun \emph{baseline}, \emph{target}, dan analisis kesenjangan.
  \item Menautkan lapisan aplikasi dengan lapisan data (akses objek data) dan
        lapisan bisnis (realisasi layanan), serta mengaitkannya dengan pola
        integrasi dan konteks Indonesia.
\end{itemize}

Tujuan pertama dan kedua mendukung CLO 3; tujuan ketiga mendukung CLO 2 dan
CLO 3; tujuan keempat mendukung CLO 4 dan CLO 6 karena mahasiswa menautkan
rancangan aplikasi dengan integrasi, API, dan konteks regulasi.

\section{Ringkasan Awal Bab}

Bab~\ref{ch:data-architecture} memodelkan \emph{data apa} yang mengalir. Bab ini
melengkapi Fase C dengan sisi \textbf{aplikasi}: \emph{perangkat lunak apa} yang
menjalankan proses bisnis dan mengolah data itu. Dalam ArchiMate, lapisan
aplikasi memuat \textbf{komponen aplikasi} (perangkat lunak yang dapat digelar),
\textbf{layanan aplikasi} (apa yang ditawarkan komponen), \textbf{fungsi
aplikasi} (perilaku internal), dan \textbf{antarmuka aplikasi} (titik akses).

Di sinilah lingkaran Fase C tertutup. Layanan aplikasi \emph{merealisasikan}
layanan bisnis dan \emph{melayani} proses bisnis; sementara fungsi aplikasi
\emph{mengakses} objek data. Jadi akses langsung ke data---yang sengaja tidak
dibuat pada lapisan bisnis (Bab~\ref{ch:data-architecture})---berada tepat di
lapisan aplikasi ini. Keluaran bab memperkuat bagian \textbf{Analisis} makalah:
peta aplikasi, matriks komponen--layanan, dan analisis kesenjangan.

\section{Pendahuluan}

Ketika organisasi berbicara tentang ``sistem'', yang dimaksud biasanya lapisan
aplikasi. Namun kesalahan umum adalah melompat ke daftar aplikasi tanpa
menautkannya ke proses bisnis dan data. Akibatnya muncul lanskap yang rumit:
banyak integrasi titik-ke-titik, duplikasi fungsi, dan ketergantungan yang sulit
ditelusuri. Arsitektur aplikasi yang baik membuat lanskap itu terbaca---komponen
apa menjalankan fungsi apa, menawarkan layanan apa, dan bergantung pada komponen
mana.

Bagi makalah, lapisan aplikasi memberi ketegasan pada rekomendasi teknologi.
Alih-alih menyebut ``bangun API'', penulis dapat menunjukkan komponen API
Gateway, layanan transfer, mesin AML, dan adaptor BI-FAST beserta relasinya.
Konsep API sebagai produk \cite{jacobson2011apis} dan tata kelola siklus hidup
API \cite{medjaoui2018continuous} memberi dasar akademik, sedangkan pembahasan
modularisasi membantu menimbang komponen mana yang dipisah atau disatukan
\cite{newman2021building}.

\section{Skenario Pembuka dan Pertanyaan Pemandu}

Setelah proses dan data transfer Bank Wanua dipetakan, tim arsitektur menghadapi
pertanyaan aplikasi: apakah membangun satu aplikasi monolitik yang menangani
semuanya, atau memisahkan gateway, layanan transfer, penyaringan AML, adaptor
BI-FAST, dan mesin rekonsiliasi menjadi komponen tersendiri? Sistem inti lama
tidak dapat diganti sekaligus, sehingga integrasi harus bertahap dan aman.

\subsection*{Pertanyaan Pemandu}
\begin{enumerate}
  \item Layanan aplikasi apa yang diperlukan untuk mewujudkan layanan bisnis
        transfer?
  \item Komponen aplikasi apa yang menjalankan fungsi-fungsi itu, dan bagaimana
        ketergantungannya?
  \item Fungsi aplikasi mana yang mengakses objek data mana?
  \item Pola integrasi apa yang sesuai dengan kondisi sistem inti dan mitra?
  \item Bagaimana kondisi lanskap saat ini berbeda dari target, dan apa
        kesenjangannya?
\end{enumerate}

\section{Mengapa Bab Ini Penting bagi Pembaca Non-TI}

Pembaca non-TI tidak perlu menulis kode, tetapi perlu memahami keputusan
aplikasi yang bersifat manajerial: membeli atau membangun, memusatkan atau
memisahkan, mengganti atau membungkus sistem lama. Keputusan ini menentukan
biaya, kecepatan perubahan, dan risiko. Dengan empat notasi sederhana---komponen,
layanan, fungsi, dan antarmuka---seorang analis dapat menjelaskan lanskap
aplikasi dan trade-off-nya tanpa tenggelam dalam detail teknis.

\section{Konsep Inti dan Glosarium Mini}

\subsection{Lapisan Aplikasi ArchiMate}

\textbf{Komponen aplikasi} (\emph{application component}) adalah unit perangkat
lunak yang dapat digelar dan diganti, misalnya API Gateway atau Mesin AML.
\textbf{Layanan aplikasi} (\emph{application service}) adalah perilaku yang
ditawarkan komponen kepada lingkungannya, misalnya Layanan API Transfer.
\textbf{Fungsi aplikasi} (\emph{application function}) adalah perilaku internal
komponen, misalnya Eksekusi Pembayaran. \textbf{Antarmuka aplikasi}
(\emph{application interface}) adalah titik akses layanan, misalnya REST API yang
patuh SNAP BI. Objek data (Bab~\ref{ch:data-architecture}) melengkapi lapisan ini
sebagai struktur pasifnya.

\subsection{Notasi Elemen Aplikasi di Archi}

Gambar~\ref{fig:app-legend} menampilkan jenis elemen lapisan aplikasi di Archi.
Semua elemen aplikasi berwarna kebiruan (biru muda); ikon di pojok kanan atas
membedakan jenisnya, dan objek data ikut ditampilkan sebagai struktur pasif
lapisan ini.

\begin{figure}[htbp]
\centering
\includegraphics[width=0.92\textwidth]{app-legend.pdf}
\caption{Jenis elemen lapisan aplikasi ArchiMate dan simbolnya di Archi.}
\label{fig:app-legend}
\end{figure}

\begin{table}[htbp]
\centering
\scriptsize
\caption{Elemen lapisan aplikasi ArchiMate, kegunaan, dan contohnya.}
\label{tab:app-elements}
\begin{tabularx}{\textwidth}{@{}p{0.24\textwidth} p{0.22\textwidth} X@{}}
\hline
\textbf{Elemen} & \textbf{Kegunaan} & \textbf{Contoh} \\
\hline
Application Component & Unit perangkat lunak yang dapat digelar/diganti. & API
Gateway, Aplikasi Transfer Dana, Mesin AML \& Limit, Adaptor BI-FAST. \\
Application Service & Perilaku yang ditawarkan komponen ke lingkungan;
merealisasikan layanan bisnis. & Layanan API Transfer, Layanan Penyaringan AML. \\
Application Function & Perilaku internal komponen; mengakses objek data. &
Otorisasi \& Validasi, Eksekusi Pembayaran, Rekonsiliasi Otomatis. \\
Application Interface & Titik akses layanan aplikasi. & REST API (SNAP BI),
Antarmuka BI-FAST. \\
Application Collaboration / Interaction & Kerja sama beberapa komponen dan
perilaku kolektifnya. & (mis.) kolaborasi gateway--mesin AML. \\
Application Process / Event & Alur perilaku aplikasi dan pemicunya. & (mis.)
alur pemrosesan permintaan, peristiwa ``permintaan masuk''. \\
Data Object & Struktur pasif---data yang diproses (Bab~6). & Data Instruksi
Transfer, Data Transaksi (Ledger). \\
\hline
\end{tabularx}
\end{table}

\subsection{Komponen, Layanan, Fungsi, dan Antarmuka}

Empat elemen ini mudah tertukar. \textbf{Layanan} menjawab ``apa yang
ditawarkan''; \textbf{fungsi} menjawab ``perilaku apa di dalam''; \textbf{komponen}
menjawab ``perangkat lunak apa yang menjalankannya''; dan \textbf{antarmuka}
menjawab ``lewat mana diakses''. Satu komponen dapat menjalankan beberapa fungsi
dan menawarkan beberapa layanan; satu layanan dapat dipakai banyak proses bisnis.

\subsection{Pola Integrasi Aplikasi}

Cara komponen saling terhubung menentukan kerumitan lanskap
(Gambar~\ref{fig:integration-patterns}). \textbf{Titik-ke-titik} menghubungkan
tiap pasang aplikasi secara langsung---sederhana di awal, tetapi cepat menjadi
kusut. \textbf{Hub/ESB} memusatkan integrasi pada satu perantara. \textbf{API-led}
menata integrasi berlapis melalui API yang dikelola sebagai produk, cocok untuk
membuka layanan ke mitra secara aman dan terkendali \cite{medjaoui2018continuous}.
Untuk rel transfer, pendekatan API-led melalui gateway sejalan dengan prinsip
``API sebagai produk''.

\begin{figure}[htbp]
\centering
\scalebox{1.0}{
\begin{tikzpicture}[
  font=\footnotesize,
  lbl/.style={font=\footnotesize\bfseries, align=center},
  app/.style={circle, draw=blue!70!black, fill=blue!10, minimum size=8mm, inner sep=0pt},
  hub/.style={rectangle, rounded corners=1mm, draw=orange!80!black, fill=orange!16,
    minimum width=13mm, minimum height=7.5mm},
  gw/.style={rectangle, rounded corners=1mm, draw=praditagreen, fill=praditagreen!14,
    minimum width=15mm, minimum height=7.5mm},
  lnk/.style={draw=gray!55}
]
  % point-to-point
  \node[lbl] at (1.4,3.6) {Titik-ke-titik};
  \foreach \a/\x/\y in {a/0/2.6, b/2.8/2.6, c/0/0.8, d/2.8/0.8, e/1.4/1.7}
    \node[app] (\a) at (\x,\y) {};
  \foreach \p/\q in {a/b,a/c,a/d,a/e,b/c,b/d,b/e,c/d,c/e,d/e} \draw[lnk] (\p)--(\q);
  % hub
  \node[lbl] at (5.9,3.6) {Hub / ESB};
  \node[hub] (h) at (5.9,1.7) {\textbf{ESB}};
  \foreach \a/\x/\y in {h1/4.5/2.6, h2/7.3/2.6, h3/4.5/0.8, h4/7.3/0.8}
    {\node[app] (\a) at (\x,\y) {}; \draw[lnk] (\a)--(h);}
  % api-led
  \node[lbl] at (10.4,3.6) {API-led};
  \node[app] (ext) at (10.4,2.6) {};
  \node[gw] (g) at (10.4,1.6) {\textbf{Gateway}};
  \foreach \a/\x in {p1/9.2, p2/10.4, p3/11.6} {\node[app] (\a) at (\x,0.8) {}; \draw[lnk] (g)--(\a);}
  \draw[lnk] (ext)--(g);
\end{tikzpicture}
}
\caption{Tiga pola integrasi aplikasi: titik-ke-titik, hub/ESB, dan API-led.}
\label{fig:integration-patterns}
\end{figure}

\subsection{TOGAF ADM Fase C: Arsitektur Aplikasi}

Fase C mencakup dua sisi: data (Bab~\ref{ch:data-architecture}) dan aplikasi.
Sisi aplikasi menyusun \emph{application landscape} baseline dan target, lalu
menganalisis kesenjangan \cite{togaf10}. Keluaran utamanya adalah katalog
komponen/layanan aplikasi, diagram kerja sama aplikasi, dan matriks komponen
terhadap fungsi/layanan bisnis. Tabel~\ref{tab:adm-phase-c-app} memetakan
langkahnya ke kasus Bank Wanua.

\begin{table}[htbp]
\centering
\scriptsize
\caption{Pemetaan TOGAF ADM Fase C (Arsitektur Aplikasi) ke artefak Bank Wanua.}
\label{tab:adm-phase-c-app}
\begin{tabularx}{\textwidth}{@{}p{0.27\textwidth} X@{}}
\hline
\textbf{Langkah Fase C (Aplikasi)} & \textbf{Artefak / keluaran pada kasus Bank Wanua} \\
\hline
Ruang lingkup dan prinsip aplikasi & Lingkup aplikasi rel transfer; prinsip
API-led, komponen modular, dan bungkus (bukan ganti) sistem inti. \\
Arsitektur aplikasi baseline & Fungsi transfer tersebar pada sistem inti dan
kanal; integrasi titik-ke-titik; AML dan rekonsiliasi manual. \\
Arsitektur aplikasi target & API Gateway, Aplikasi Transfer Dana, Mesin AML,
Adaptor BI-FAST, dan Mesin Rekonsiliasi sebagai komponen; layanan aplikasi yang
mewujudkan layanan bisnis. \\
Analisis kesenjangan & Selisih baseline--target: komponen yang dibangun,
diintegrasikan, atau dibungkus dari sistem lama. \\
Masukan ke fase berikutnya & Kebutuhan lapisan teknologi: node, platform,
jaringan, dan sertifikasi untuk menggelar komponen (bab berikutnya). \\
\hline
\end{tabularx}
\end{table}

\subsection{Contoh Interaksi Antar-Elemen Aplikasi}

Gambar~\ref{fig:app-contoh} memakai contoh restoran cepat saji untuk menunjukkan
interaksi antar-elemen aplikasi: komponen \emph{Aplikasi Kasir (POS)}
menjalankan fungsi \emph{Pemrosesan Pesanan}, yang mewujudkan \emph{Layanan
Pemesanan} (merealisasikan layanan bisnis Layanan Makan di Tempat dan melayani
proses Menyajikan Pesanan) serta mengakses objek data \emph{Data Pesanan}.

\begin{figure}[htbp]
\centering
\includegraphics[width=\textwidth,height=0.42\textheight,keepaspectratio]{app-contoh.pdf}
\caption{Contoh interaksi antar-elemen aplikasi pada kasus restoran cepat saji.}
\label{fig:app-contoh}
\end{figure}

\subsection*{Glosarium Mini}
\begin{description}
  \item[Komponen aplikasi] Unit perangkat lunak yang dapat digelar, misalnya API
        Gateway.
  \item[Layanan aplikasi] Perilaku yang ditawarkan komponen; merealisasikan
        layanan bisnis.
  \item[Fungsi aplikasi] Perilaku internal komponen yang mengakses objek data.
  \item[Antarmuka aplikasi] Titik akses layanan, misalnya REST API.
  \item[Application landscape] Peta komponen aplikasi dan ketergantungannya.
  \item[Pola integrasi] Cara komponen terhubung (titik-ke-titik, hub/ESB,
        API-led).
\end{description}

\section{Kerangka Berpikir}

Alur berpikir bab ini bergerak dari layanan bisnis menuju lanskap aplikasi
(Gambar~\ref{fig:app-flow}). Layanan bisnis menentukan layanan aplikasi yang
harus ada; layanan aplikasi diwujudkan oleh fungsi dan komponen; komponen
saling bergantung membentuk lanskap; fungsi mengakses objek data; dan pola
integrasi menata semuanya. Hasilnya menjadi bahan analisis dan pratinjau
kebutuhan teknologi.

\begin{figure}[htbp]
\centering
\scalebox{0.86}{
\begin{tikzpicture}[
  box/.style={rectangle, rounded corners=1.5mm, draw=blue!70!black,
    fill=blue!8, thick, align=center, text width=30mm, minimum height=11mm},
  outnode/.style={rectangle, rounded corners=1.5mm, draw=orange!80!black,
    fill=orange!16, thick, align=center, text width=32mm, minimum height=11mm},
  link/.style={-Latex, thick, draw=blue!70!black}
]
  \node[box] (bs) {Layanan bisnis\\(dari Bab 5)};
  \node[box, right=7mm of bs] (as) {Layanan \&\\komponen aplikasi};
  \node[box, right=7mm of as] (land) {Lanskap \&\\ketergantungan};
  \node[box, below=8mm of land] (data) {Akses objek data\\(Bab 6)};
  \node[outnode, left=7mm of data] (integ) {Pola integrasi\\\& kesenjangan};
  \node[outnode, left=7mm of integ] (an) {Analisis \&\\pratinjau teknologi};
  \draw[link] (bs) -- (as); \draw[link] (as) -- (land); \draw[link] (land) -- (data);
  \draw[link] (data) -- (integ); \draw[link] (integ) -- (an);
  \draw[link, dashed] (an) -- (bs);
\end{tikzpicture}
}
\caption{Alur berpikir lapisan aplikasi untuk makalah publikasi.}
\label{fig:app-flow}
\end{figure}

\section{Konteks Indonesia}

Arsitektur aplikasi layanan keuangan Indonesia dibentuk oleh standar dan
regulasi. Standar Open API pembayaran nasional membakukan antarmuka dan format
data yang dibuka ke mitra \cite{snapbi}; karena itu antarmuka aplikasi pada rel
transfer dirancang patuh SNAP BI sejak awal. Penyelenggara sistem elektronik
wajib mengelola sistem, keandalan, dan jejaknya \cite{pp712019}, sedangkan
pemrosesan data pribadi pada komponen aplikasi harus memenuhi perlindungan data
pribadi \cite{uupdp2022}. Realitas sistem inti lama membuat pendekatan
\emph{bungkus bertahap} (adaptor + gateway) lebih realistis daripada penggantian
sekaligus.

\section{Studi Kasus Utama: Arsitektur Aplikasi Transfer Bank Wanua}

Layanan dan proses bisnis pada Bab~\ref{ch:business-architecture} diwujudkan oleh
sekumpulan komponen aplikasi pada Gambar~\ref{fig:app-transfer}.

\begin{figure}[htbp]
\centering
\includegraphics[width=\textwidth,height=0.62\textheight,keepaspectratio]{app-transfer.pdf}
\caption{Lanskap aplikasi (ArchiMate) API Transfer Dana Bank Wanua: komponen
menjalankan fungsi, fungsi merealisasikan layanan aplikasi, antarmuka mengekspos
layanan, dan komponen saling melayani.}
\label{fig:app-transfer}
\end{figure}

\subsection{Cara Membaca Diagram Aplikasi}

Gambar~\ref{fig:app-transfer} dibaca dari kiri ke kanan. Kolom kiri memuat
\textbf{komponen aplikasi}: API Gateway (pintu masuk yang aman dan patuh SNAP
BI), Aplikasi Transfer Dana (orkestrator), Mesin AML \& Limit, Adaptor BI-FAST,
Adaptor Core Banking, dan Mesin Rekonsiliasi \& Notifikasi. Tiap komponen
\emph{menjalankan} (assignment) sebuah \textbf{fungsi aplikasi} di kolom tengah:
Otorisasi \& Validasi, Penyaringan AML \& Limit, Eksekusi Pembayaran,
Rekonsiliasi Otomatis, dan Notifikasi. Tiap fungsi lalu \emph{merealisasikan}
sebuah \textbf{layanan aplikasi} di kolom kanan (Layanan API Transfer, Penyaringan
AML, Rekonsiliasi, dan Notifikasi). Antarmuka REST API \emph{mengekspos} layanan
transfer, dan komponen-komponen saling \emph{melayani} (serving) untuk membentuk
alur: gateway melayani orkestrator, sedangkan mesin AML, adaptor BI-FAST, dan
mesin rekonsiliasi melayani orkestrator; adaptor core banking melayani adaptor
BI-FAST.

\subsection{Penyelarasan Bisnis, Aplikasi, dan Data}

Kekuatan lapisan aplikasi terlihat saat ketiga lapisan Fase B--C disandingkan
(Gambar~\ref{fig:app-stack}). Diagram ini menutup rantai keterlacakan
end-to-end: \textbf{layanan aplikasi merealisasikan layanan bisnis} dan
\textbf{melayani proses bisnis}; \textbf{fungsi aplikasi merealisasikan layanan
aplikasi}; \textbf{komponen menjalankan fungsi}; dan \textbf{fungsi aplikasi
mengakses objek data}.

\begin{figure}[htbp]
\centering
\includegraphics[width=\textwidth,height=0.62\textheight,keepaspectratio]{app-stack.pdf}
\caption{Penyelarasan lintas-lapisan (ArchiMate) untuk rel transfer: bisnis,
aplikasi (layanan dan komponen/fungsi), serta data. Fungsi aplikasi mengakses
objek data---di sinilah akses data langsung berada, bukan pada proses bisnis.}
\label{fig:app-stack}
\end{figure}

Inilah jawaban atas catatan pada Bab~\ref{ch:data-architecture}: proses bisnis
mengakses objek bisnis, sedangkan \emph{objek data diakses langsung oleh fungsi
aplikasi}. Dengan rantai lengkap---proses bisnis dilayani layanan aplikasi, yang
diwujudkan fungsi, yang mengakses data---setiap kebutuhan dapat ditelusuri dari
motivasi sampai ke data dan komponen yang mewujudkannya.

\section{Langkah Analisis dan Artefak Pendukung}

\subsection{Langkah 1: Ambil Layanan dan Proses Bisnis}
Mulai dari layanan dan proses bisnis (Bab~\ref{ch:business-architecture}) yang
perlu didukung aplikasi.

\subsection{Langkah 2: Tentukan Layanan Aplikasi}
Untuk tiap layanan bisnis, tentukan layanan aplikasi yang merealisasikannya dan
proses bisnis yang dilayaninya.

\subsection{Langkah 3: Rancang Komponen dan Fungsi}
Tentukan komponen yang menjalankan fungsi dan mewujudkan layanan aplikasi.
Pertimbangkan memisahkan atau menyatukan komponen sesuai kebutuhan perubahan.

\subsection{Langkah 4: Tautkan Fungsi ke Objek Data}
Nyatakan relasi akses baca/tulis antara fungsi aplikasi dan objek data
(Bab~\ref{ch:data-architecture}).

\subsection{Langkah 5: Pilih Pola Integrasi}
Tetapkan pola integrasi (titik-ke-titik, hub/ESB, API-led) dan jelaskan
alasannya terhadap kondisi sistem inti dan mitra.

\subsection{Langkah 6: Susun Baseline dan Target}
Bandingkan lanskap aplikasi saat ini dan target, lalu catat kesenjangan.

\subsection{Langkah 7: Kaitkan dengan Standar dan Regulasi}
Tautkan antarmuka dengan SNAP BI, dan pemrosesan data pribadi dengan UU PDP serta
PP 71.

\subsection{Langkah 8: Tulis Narasi Analisis Aplikasi}
Ubah lanskap dan matriks menjadi paragraf yang menjelaskan keputusan aplikasi dan
implikasinya bagi lapisan teknologi.

\begin{table}[htbp]
\centering
\scriptsize
\caption{Artefak pendukung Bab~7 dan fungsinya dalam makalah.}
\label{tab:artefak-bab6-app}
\begin{tabularx}{\textwidth}{@{}p{0.26\textwidth} p{0.30\textwidth} X@{}}
\hline
\textbf{Artefak} & \textbf{Isi minimum} & \textbf{Dipakai untuk bagian makalah} \\
\hline
Katalog komponen/layanan & Daftar komponen, layanan, fungsi & Analisis domain
aplikasi. \\
Diagram lanskap aplikasi & Komponen \& ketergantungannya & Analisis dan lampiran. \\
Diagram penyelarasan lintas-lapisan & Bisnis--aplikasi--data & Analisis
keterlacakan end-to-end. \\
Analisis pola integrasi & Pilihan pola + alasan & Diskusi \& implikasi
manajerial. \\
\hline
\end{tabularx}
\end{table}

\section{Integrasi ke Makalah Publikasi}

Lapisan aplikasi memperkuat bagian \textbf{Analisis} dan \textbf{Diskusi}.
Diagram penyelarasan menjadi bukti keterlacakan dari kebutuhan bisnis sampai ke
komponen dan data. Contoh paragraf:

\begin{quote}
Analisis aplikasi menunjukkan bahwa layanan transfer real-time paling baik
diwujudkan oleh lanskap API-led: sebuah API Gateway yang patuh SNAP BI, aplikasi
transfer sebagai orkestrator, serta komponen AML, adaptor BI-FAST, dan mesin
rekonsiliasi yang terpisah. Pemisahan ini memungkinkan sistem inti dibungkus
bertahap tanpa penggantian sekaligus, sekaligus menjaga keterlacakan akses data
pada tiap fungsi.
\end{quote}

\section{Diskusi Kritis, Trade-off, dan Perangkap Umum}

\begin{table}[htbp]
\centering
\scriptsize
\caption{Trade-off dalam arsitektur aplikasi.}
\label{tab:tradeoff-bab6-app}
\begin{tabularx}{\textwidth}{@{}p{0.24\textwidth} X X@{}}
\hline
\textbf{Keputusan} & \textbf{Jika terlalu ringan} & \textbf{Jika terlalu berat} \\
\hline
Granularitas komponen & Monolit sulit diubah & Terlalu banyak komponen, integrasi
rumit. \\
Pola integrasi & Titik-ke-titik menjadi kusut & Platform integrasi mahal dan
berat. \\
Sistem inti lama & Diabaikan, integrasi rapuh & Diganti sekaligus, risiko tinggi. \\
Keamanan API & Antarmuka terbuka tanpa kendali & Kontrol berlebihan menghambat
adopsi mitra. \\
\hline
\end{tabularx}
\end{table}

Perangkap umum: (1) mendaftar aplikasi tanpa menautkannya ke proses dan data;
(2) integrasi titik-ke-titik yang menumpuk; (3) mengganti sistem inti sekaligus;
dan (4) membuka API tanpa tata kelola dan kepatuhan standar.

\section{Latihan Terapan Individual}

\subsection*{Latihan~7.1: Lanskap Aplikasi}
Untuk organisasi studi kasus Anda, susun lanskap aplikasi: komponen, fungsi,
layanan, dan antarmuka, beserta ketergantungannya.

\subsection*{Latihan~7.2: Penyelarasan Lintas-Lapisan}
Buat diagram yang menautkan layanan bisnis, layanan/komponen aplikasi, dan objek
data untuk satu kapabilitas sasaran.

\subsection*{Latihan~7.3: Pola Integrasi dan Baseline--Target}
Pilih pola integrasi beserta alasannya, lalu bandingkan lanskap baseline dan
target dan catat kesenjangan.

\subsection*{Latihan~7.4: Paragraf Analisis Aplikasi}
Tulis satu paragraf analisis aplikasi yang menautkan keputusan komponen dengan
kebutuhan bisnis dan implikasi teknologi.

\begin{table}[htbp]
\centering
\scriptsize
\caption{Rubrik latihan Bab~7.}
\label{tab:rubrik-bab6-app}
\begin{tabularx}{\textwidth}{@{}p{0.25\textwidth} X X@{}}
\hline
\textbf{Komponen} & \textbf{Indikator baik} & \textbf{Perlu diperbaiki bila} \\
\hline
Lanskap aplikasi & Komponen, layanan, fungsi jelas \& tertaut & Sekadar daftar
aplikasi tanpa relasi. \\
Penyelarasan lintas-lapisan & Menautkan bisnis--aplikasi--data & Aplikasi berdiri
sendiri tanpa jejak. \\
Pola integrasi & Pilihan beralasan terhadap konteks & Tanpa pertimbangan sistem
lama/mitra. \\
Narasi analisis & Menautkan komponen ke kebutuhan \& implikasi & Deskriptif
tanpa argumen. \\
\hline
\end{tabularx}
\end{table}

\section{Ringkasan, Refleksi, dan Jembatan}

\subsection*{Poin Kunci}
\begin{itemize}
  \item Lapisan aplikasi memuat komponen, layanan, fungsi, dan antarmuka.
  \item Layanan aplikasi \emph{merealisasikan} layanan bisnis dan \emph{melayani}
        proses bisnis; fungsi aplikasi \emph{mengakses} objek data.
  \item Akses data langsung berada di lapisan aplikasi, bukan pada proses bisnis.
  \item Pola integrasi (titik-ke-titik, hub/ESB, API-led) menentukan kerumitan
        lanskap.
  \item Fase C (Aplikasi) menghasilkan baseline, target, dan analisis kesenjangan.
\end{itemize}

\subsection*{Uji Mandiri}
\begin{enumerate}
  \item Jelaskan perbedaan komponen, layanan, fungsi, dan antarmuka aplikasi.
  \item Mengapa akses objek data ditempatkan pada fungsi aplikasi, bukan pada
        proses bisnis?
  \item Ambil satu layanan bisnis pada kasus Anda dan sebutkan layanan serta
        komponen aplikasi yang mewujudkannya.
\end{enumerate}

\subsection*{Jembatan ke Bab Berikutnya}
Bab berikutnya (Bab~\ref{ch:technology-architecture}) turun ke \textbf{lapisan
teknologi}: node, perangkat, jaringan,
dan \emph{system software} yang menggelar komponen aplikasi ini, serta artifact
yang mewujudkan objek data. Katalog komponen dan diagram penyelarasan menjadi
masukan langsung untuk memilih platform, pola penggelaran, dan arsitektur
teknologi acuan.

\section{Bacaan Lanjutan dan Lampiran}

\subsection*{Bacaan Inti}
\begin{itemize}
  \item The Open Group untuk Fase C ADM dan artefak arsitektur aplikasi
        \cite{togaf10}.
  \item Newman untuk modularisasi, komponen, dan pola integrasi
        \cite{newman2021building}.
\end{itemize}

\subsection*{Bacaan untuk Memperdalam}
\begin{itemize}
  \item Untuk API sebagai produk dan tata kelola siklus hidupnya
        \cite{jacobson2011apis,medjaoui2018continuous}.
  \item Untuk konteks Indonesia: standar API pembayaran \cite{snapbi},
        penyelenggaraan sistem elektronik \cite{pp712019}, dan perlindungan data
        pribadi \cite{uupdp2022}.
\end{itemize}

\subsection*{Lampiran Kerja Bab~7}
\begin{description}
  \item[Katalog komponen/layanan] Daftar komponen, fungsi, layanan, antarmuka.
  \item[Diagram lanskap aplikasi] Komponen dan ketergantungannya.
  \item[Diagram penyelarasan] Bisnis--aplikasi--data.
  \item[Analisis integrasi] Pola integrasi dan kesenjangan baseline--target.
\end{description}
\setlength{\emergencystretch}{0pt}
